% ----------------------------------------------------------
\chapter{Fundamentação Teórica}\label{cap:fundamentacao}
% ----------------------------------------------------------
Este capítulo reúne os conceitos e as diretrizes que fundamentam o desenvolvimento do trabalho, organizados em nove eixos: o Diabetes Mellitus tipo 2 e o desafio na saúde pública; o rastreamento, a APS e as diretrizes clínicas; a Pesquisa Nacional de Saúde de 2013 e os determinantes da saúde; as variáveis antropométricas e clínicas como classificadoras de DM2; a mineração de dados; o pré-processamento e o tratamento de dados; os algoritmos de \textit{Machine Learning}; a avaliação de modelos classificatórios; a Inteligência Artificial Explicável; e a multicolinearidade e seus efeitos sobre a modelagem e a interpretabilidade. Cada eixo dialoga com uma decisão apresentada na Metodologia, de modo que o referencial justifique as etapas do \textit{pipeline} de desenvolvimento do modelo classificatório, em vez de reunir definições isoladas.

\section{O Diabetes Mellitus Tipo 2 e o Desafio na Saúde Pública} \label{sec:dm2-desafio}

O Diabetes Mellitus tipo 2 (DM2) é uma das principais doenças epidêmicas do mundo, representando cerca de 90\% dos casos de diabetes ao redor do globo \cite{who2024diabetes}. Conforme apontado pela literatura especializada em mecanismos patogênicos \cite{mlynarska2025type}, o DM2 é caracterizado por uma progressão silenciosa, na qual a resistência à insulina e o declínio funcional das células beta-pancreáticas\footnote{Células responsáveis pela produção, armazenamento e secreção do hormônio insulina.} antecedem o diagnóstico clínico. Durante anos, ocorre a convivência assintomática com a hiperglicemia\footnote{Condição caracterizada por níveis excessivamente altos de glicose (açúcar) no sangue.}, configurando uma janela temporal propícia a danos celulares progressivos. Essa característica silenciosa colabora para que o diagnóstico clínico ocorra de forma tardia, etapa na qual complicações como nefropatia, retinopatia, neuropatia e doença cardiovascular já podem estar presentes \cite{mlynarska2025type}. 

De acordo com a 11ª edição do Atlas de Diabetes da Federação Internacional de Diabetes \cite{IDF2025}, o Brasil ocupava, em 2024, o 6º lugar no ranking global das dez nações com o maior número de adultos (entre 20 e 79 anos) vivendo com a doença. Além da alta prevalência, o país encontra-se na 9ª posição mundial em volume de pacientes com diabetes não diagnosticado e é classificado como a 3ª nação com o maior gasto total em saúde (em dólares americanos) direcionado ao tratamento da DM. Segundo \textcite{bezerra2025morbidade} o custo médio global de internações de pacientes por Diabetes Mellitus (DM) no Sistema Único de Saúde (SUS), entre janeiro de 2014 e dezembro de 2024, foi de R\$ 832,93, o que evidencia que as consequências geradas pelo DM transcendem a morbidade individual e atingem diretamente a viabilidade financeira do SUS. 

Nesse contexto, conforme as diretrizes atualizadas da Sociedade Brasileira de Diabetes \cite{SBD2025} a triagem precoce configura-se como estratégia de custo-efetividade potencial, ao iniciar o contato com o indivíduo em risco para uma fase anterior ao início de complicações clínicas de alto custo assistencial, já que o diagnóstico precoce e o tratamento nas fases iniciais são fundamentais para a prevenção da evolução para complicações crônicas e a consequente necessidade de encaminhamento para a atenção especializada \cite{Rodacki2025}.

A triagem não se constitui, de forma isolada, como instrumento suficiente para a redução de internações, visto que o desfecho clínico depende de diversos fatores sociais, de acesso à atenção à saúde e do comprometimento individual ao tratamento \cite{geremia2020atenccao}. O que a triagem viabiliza é o encaminhamento assertivo de indivíduos para a confirmação diagnóstica laboratorial, transferindo o foco da intervenção reativa para a prevenção orientada por evidência \cite{brasil2013}. Para fins de parâmetro, o custo unitário do exame de Dosagem de Hemoglobina Glicosilada (HbA1c) tabelado no Sistema de Gerenciamento da Tabela de Procedimentos, Medicamentos e OPM do SUS (SIGTAP) é de R\$ 7,86 \cite{sigtap_hemoglobina_2026}. Embora este valor isoladamente não determine ganhos clínicos ou financeiros, o mesmo ilustra a discrepância de custos entre a fase de rastreamento e a fase de manejo de complicações instaladas \cite{bezerra2025morbidade}, reforçando a relevância econômica do investimento em tecnologias de triagem direcionada na Atenção Primária à Saúde.

\subsection{Rastreamento, APS e Diretrizes Clínicas}\label{sec:rastreamento-aps}
A Sociedade Brasileira de Diabetes publicou, em 2025, uma atualização significativa de suas diretrizes de rastreamento e diagnóstico, cujo marco central é o rebaixamento da idade de início do rastreamento universal de 45 para 35 anos \cite{Rodacki2025}. Essa mudança, baseada em evidências epidemiológicas que apontam para a crescente incidência do DM2 em adultos jovens, representa um avanço substancial na política pública de prevenção, de maneira em que a detecção precoce passa a abranger uma faixa etária anteriormente excluída dos protocolos sistemáticos de triagem \cite{SBD2025}. Adicionalmente, a edição de 2025 incorpora, pela primeira vez, um capítulo dedicado ao uso de ferramentas digitais e Inteligência Artificial no manejo do diabetes, legitimando a adoção de modelos computacionais como suporte à decisão clínica \cite{SBD2025}. Diante disso, o presente trabalho encontra-se diretamente alinhado à direção que a própria sociedade médica de referência nacional aponta como caminho para o aprimoramento do cuidado ao paciente diabético.

Do ponto de vista epidemiológico, a distinção entre triagem populacional e diagnóstico clínico é central para a compreensão dos objetivos e das métricas de desempenho que devem orientar instrumentos de rastreamento \cite{Swaminathan2024, Rodacki2025}. A triagem tem como premissa fundamental a maximização da Sensibilidade (\textit{Recall}), métrica que expressa a capacidade do instrumento de classificar corretamente indivíduos com perfil de risco metabólico associado ao Diabetes Mellitus tipo 2, admitindo-se uma taxa controlada de falsos positivos como custo operacional aceitável, já que o erro de triagem (falso negativo) permite que o paciente permaneça na fase inicial da doença por uma janela de tempo prolongada na qual a hiperglicemia instaura danos celulares crônicos e silenciosos, conforme apontado por \textcite{mlynarska2025type}. O diagnóstico clínico formal exige alta especificidade, priorizando a correta exclusão de indivíduos saudáveis a fim de evitar intervenções terapêuticas desnecessárias \cite{mlynarska2025type, ADA2010, SBD2025}. Conforme apontado por \textcite{Rodacki2025}, o rastreamento torna-se mais eficaz quanto maior a prevalência da doença na população-alvo, de forma com que a seleção de indivíduos com maior probabilidade de resultado positivo visa maximizar a efetividade da iniciativa e a alocação racional dos recursos disponíveis na rede pública de saúde.

No que se refere às ferramentas de triagem atualmente disponíveis para o contexto brasileiro, o \textit{Finnish Diabetes Risk Score} (FINDRISC), adaptado e validado para o português brasileiro por \textcite{conceiccao2020translation}, configura-se como um dos instrumentos mais adotados na Atenção Primária. O escore opera por meio da atribuição de pesos fixos a variáveis como idade, índice de massa corporal, circunferência abdominal, prática de atividade física e histórico familiar de diabetes, gerando uma pontuação ordinal que estratifica o indivíduo em faixas de risco \cite{conceiccao2020translation}. Conforme apontado pela revisão sistemática de \textcite{CarrilloLarco2020} o FINDRISC foi originalmente calibrado em populações europeias, de maneira em que sua aplicação direta a uma população geneticamente heterogênea como a brasileira pode introduzir vieses sistemáticos de classificação. A estrutura de pesos estáticos, embora válida, não permite a captura de interações não lineares entre fatores de risco que, em combinação, elevam expressivamente a probabilidade de hiperglicemia crônica em perfis populacionais específicos, como os documentados nas regiões Norte e Nordeste do Brasil \cite{muzy2021prevalencia}.

O \textit{BrDMrisc}, escore de risco para diabetes incidente derivado da coorte ELSA-Brasil \cite{bracco2026brdmrisc}, representa um avanço metodológico relevante em relação ao FINDRISC ao ser desenvolvido e calibrado sobre dados populacionais brasileiros. Contudo, o instrumento também opera sobre pesos fixos estimados em contexto longitudinal prospectivo, abordagem que, embora metodologicamente robusta, não captura a variabilidade de risco decorrente das interações complexas entre determinantes sociais, antropométricos e clínicos presentes em uma população de heterogeneidade continental como a brasileira \cite{islam2024comprehensive, muzy2021prevalencia}.

A viabilidade de incorporação de modelos classificatórios ao contexto da APS encontra respaldo na capacidade de algoritmos de \textit{Machine Learning} (ML) em capturar padrões complexos em grandes conjuntos de dados, o que inclui interações não lineares entre fatores de risco. Além disso, a crescente digitalização do sistema de saúde brasileiro juntamente da adoção do Prontuário Eletrônico do Cidadão (PEC) no âmbito do Sistema de Informação em Saúde para a Atenção Básica (SISAB) já integra o gerenciamento do cuidado das equipes multidisciplinares \cite{brasil_esus_aps}, configurando uma infraestrutura de captura de dados clínicos rotineiros que, até o presente momento, opera predominantemente como repositório administrativo e não como fonte ativa de inteligência epidemiológica. Contudo, a literatura aponta que o modelo assistencial para condições crônicas ainda opera, de forma predominante, fundamentado na demanda espontânea do paciente \cite{geremia2020atenccao}, de maneira em que indivíduos assintomáticos em fase de risco metabólico elevado permanecem subnotificados até a manifestação clínica de complicações. A aplicação de algoritmos de ML sobre as variáveis rotineiras capturadas pelo PEC apresenta o potencial de romper com essa lógica reativa, transformando o volume de informações clínicas já disponíveis em inteligência epidemiológica ativa, a fim de viabilizar a triagem, classificação e encaminhamento assertivo de indivíduos em fases iniciais para a confirmação diagnóstica laboratorial. A resistência profissional à adoção tecnológica na APS, contudo, figura entre os principais desafios documentados na literatura \cite{geremia2020atenccao}, de forma com que a aceitabilidade clínica de modelos classificatórios depende, em medida expressiva, da capacidade do sistema de apresentar explicações compreensíveis ao profissional de saúde sobre os fatores que motivaram cada classificação. Diante disso, a necessidade de acoplar modelos classificatórios a interfaces de Inteligência Artificial Explicável (XAI) configura-se não apenas como escolha técnica, mas como requisito de adoção clínica \cite{salih2025perspective}.


\section{A Pesquisa Nacional de Saúde (PNS 2013) e os Determinantes da Saúde} \label{sec:pns}

A Pesquisa Nacional de Saúde (PNS) de 2013 é um dos maiores inquéritos populacionais em termos de abrangência e complexidade metodológica já conduzidos no Brasil e integra o Sistema Integrado de Pesquisas Domiciliares (SIPD) do Instituto Brasileiro de Geografia e Estatística \cite{IBGE-SIPD}. O plano amostral adotado pela pesquisa baseia-se na amostragem por conglomerado\footnote{Método probabilístico em que a população é dividida em grupos naturais ou geográficos.} em três estágios de seleção, com estratificação das unidades primárias de amostragem \cite{IBGE_PNS2013_Exames}. No primeiro estágio, foram selecionados os setores censitários ou composições de setores; no segundo estágio, elegeu-se um número fixo de domicílios particulares permanentes em cada unidade primária; e no terceiro estágio, realizou-se a seleção de um único morador adulto por domicílio para a aplicação do questionário individual de saúde. Essa arquitetura amostral complexa acarreta na utilização de pesos amostrais\footnote{Valores numéricos atribuídos a cada indivíduo ou elemento de uma pesquisa que indicam quantas pessoas da população geral aquela observação representa.} nos procedimentos analíticos, a fim de garantir que as estimativas produzidas reflitam a representatividade da população adulta brasileira, e não apenas a distribuição observada na amostra bruta \cite{IBGE_PNS2013_Exames}.

Um diferencial metodológico da edição de 2013 reside na presença de uma Subamostra Laboratorial, na qual foram coletados exames de sangue e urina de 8.952 participantes adultos, incluindo a Dosagem de Hemoglobina Glicosilada (HbA1c) \cite{IBGE_PNS2013_Exames}. A disponibilidade desse marcador laboratorial viabiliza a identificação de indivíduos com hiperglicemia crônica clinicamente significativa independentemente do autorrelato diagnóstico, de modo que o desfecho do presente trabalho (HbA1c $\geq$ 6,5\%, conforme critério da American Diabetes Association \cite{ADA2010}) fundamenta-se em evidência clínica objetiva e não em percepção subjetiva de saúde. A importância dessa distinção é evidenciada pelo fenômeno da subnotificação diagnóstica, já que enquanto a prevalência autorreferida de diabetes na PNS 2013 foi de 6,2\%, a prevalência corrigida por critério laboratorial alcançou 9,4\%, implicando uma proporção nacional de subnotificação de 42,5\%, que atingiu 72,8\% na Região Norte do país \cite{muzy2021prevalencia}. Esse conjunto de evidências demonstra que o autorrelato isolado configura-se como instrumento insuficiente para a vigilância epidemiológica e treinamento de modelos de ML classificatórios para identificação de risco metabólico associado ao DM2, corroborado pela literatura nacional consolidada \cite{iser2015prevalencia, reis2022evolution}.

Em virtude da elevada taxa de não resposta na etapa de coleta biológica, estimada em aproximadamente 52\% em relação à subamostra originalmente prevista, o IBGE desenvolveu e publicou, em 2020, uma Nota Técnica \cite{IBGE} dedicada ao recálculo do peso amostral laboratorial, variável denominada \textit{Peso\_Lab} no arquivo dos microdados disponibilizado pelo IBGE \cite{IBGE_PNS2013_Exames}. Esse peso recalibrado incorpora correções para as distorções provenientes do desenho amostral complexo e das perdas na coleta de material biológico, de forma que as taxas, médias e proporções estimadas a partir da subamostra reflitam o perfil epidemiológico ajustado da população adulta brasileira. Conforme orientação metodológica do IBGE, o Peso\_Lab normalizado para o tamanho da subamostra (n = 8.952) permite a estimação de proporções e taxas representativas, mas não a projeção de volumes populacionais absolutos \cite{IBGE_PNS2013_Exames}.

A diversidade da população brasileira, evidenciada pelas disparidades regionais na prevalência e subnotificação do DM2 documentadas por \textcite{muzy2021prevalencia, iser2015prevalencia, reis2022evolution}, não se explica exclusivamente por fatores biológicos e antropométricos. A literatura especializada em epidemiologia social aponta que variáveis sociodemográficas como escolaridade, raça/cor e macrorregião geográfica atuam como indicadores aproximados dos determinantes sociais da saúde, refletindo iniquidades históricas no acesso à atenção primária, no letramento em saúde e na capacidade de adesão a condutas preventivas \cite{deAlmeida-Pititto09112015}. A escolaridade, em particular, condiciona a capacidade do indivíduo de compreender sinais de alerta, acessar serviços de saúde de forma proativa e aderir a mudanças de estilo de vida recomendadas, de maneira que sua inclusão em modelos classificatórios não visa apenas capturar risco metabólico intrínseco, mas também a dimensão de vulnerabilidade diagnóstica estrutural. A variável raça/cor, por sua vez, captura iniquidades históricas e barreiras de acesso diferencial ao sistema de saúde, além de estresse alostático crônico associado a contextos de discriminação sistêmica, fatores que a literatura epidemiológica brasileira associa a desfechos metabólicos adversos \cite{deAlmeida-Pititto09112015, muzy2021prevalencia}.

Por fim, uma limitação própria ao desenho transversal da PNS 2013 merece fundamentação teórica explícita, visto que ela impõe restrições à inclusão de variáveis comportamentais no modelo classificatório proposto. Em estudos transversais, a exposição e o desfecho são mensurados simultaneamente, impossibilitando estabelecer precedência temporal entre causa e efeito, conforme apontado por \textcite{rothman2008modern}. Diante disso, variáveis comportamentais como prática de exercício físico e padrão alimentar estão sujeitas ao fenômeno de causalidade reversa, na qual o comportamento observado no momento da coleta pode ser consequência do diagnóstico clínico preexistente e não um determinante anterior ao desfecho \cite{rothman2008modern,dimova2020diagnosis}. Como evidenciado por \textcite{dimova2020diagnosis}, o diagnóstico de DM2 frequentemente configura-se como um momento de mudança comportamental induzida pela orientação clínica, na qual o paciente recém-diagnosticado passa a adotar hábitos mais saudáveis por prescrição médica, assim um modelo treinado sobre dados transversais poderia aprender a associar exercício físico e alimentação adequada à presença da doença crônica, invertendo a relação causal real. Diante disso, a exclusão deliberada das variáveis comportamentais do conjunto de classificadores do presente trabalho trata-se de uma decisão metodológica fundamentada na teoria epidemiológica consolidada \cite{rothman2008modern}, a fim de preservar a validade interna das associações aprendidas pelo algoritmo.

\section{Variáveis Antropométricas e Clínicas como Classificadoras de DM2}\label{sec:variaveis}

A seleção das variáveis classificadoras voltadas à identificação de risco metabólico associado ao DM2 fundamentou-se em evidências clínicas e epidemiológicas consolidadas, de modo que cada variável incorporada ao conjunto de classificadores apresente justificativa biológica e empírica para a sua associação com o desfecho de hiperglicemia crônica. Neste trabalho, as variáveis antropométricas constituem o núcleo do bloco biológico do modelo e, de acordo com \textcite{mlynarska2025type}, foram selecionadas em razão de sua capacidade de capturar dimensões complementares do fenótipo individual, em especial aquelas associadas à adiposidade visceral\footnote{Acúmulo de gordura intra-abdominal ao redor de órgãos vitais como fígado, pâncreas e intestinos.} e à resistência à insulina\footnote{Condição em que o corpo não responde adequadamente à insulina, hormônio que facilita a entrada de glicose nas células.}.

O Índice de Massa Corporal (IMC), calculado pela razão entre o peso corporal em quilogramas e o quadrado da altura em metros, fórmula clássica de Quetelet \cite{eknoyan2008adolphe}, constitui o preditor basal de sobrepeso e obesidade mais amplamente utilizado em estudos epidemiológicos e diretrizes clínicas, em virtude de sua simplicidade e da disponibilidade de pontos de corte internacionalmente consolidados \cite{mlynarska2025type}. Entretanto, a literatura de \textcite{Feller2010} evidencia que o IMC apresenta limitações relevantes para a avaliação do risco metabólico, visto que a medida não distingue adiposidade visceral de adiposidade subcutânea\footnote{Camada de gordura localizada logo abaixo da pele.}, compartimentos adiposos com propriedades metabólicas substancialmente distintas \cite{Feller2010}. Os textos de \textcite{Feller2010,mlynarska2025type} relatam que o tecido adiposo visceral, ao contrário do subcutâneo, é metabolicamente ativo e secreta adipocinas inflamatórias que induzem resistência à insulina e disfunção endotelial\footnote{Alteração precoce na camada que reveste o interior dos vasos sanguíneos (endotélio).}, configurando-se como determinante patogênico do DM2.

Diante dessa limitação, a Circunferência da Cintura (CC) emerge como preditor superior para a quantificação da adiposidade visceral. O estudo de \textcite{Feller2010}, conduzido com 9.753 homens e 15.491 mulheres participantes da coorte EPIC-Potsdam\footnote{Projeto populacional de coorte prospectiva na Alemanha.}, demonstrou que um incremento de 1 cm na circunferência da cintura está associado a um aumento de 8\% no risco relativo de desenvolvimento de DM2, tanto em homens quanto em mulheres \cite{Feller2010}. De forma ainda mais expressiva, o estudo evidenciou que indivíduos com IMC inferior a 25 kg/m² mas com circunferência da cintura elevada apresentaram risco relativo de DM2 de 3,62 vezes em homens e 2,74 vezes em mulheres em comparação ao grupo de referência, risco equivalente ou superior ao observado em indivíduos com pré-obesidade e cintura normal \cite{Feller2010}. Esse achado evidencia que a utilização isolada do IMC como critério de triagem configura-se insuficiente para a identificação de perfis de risco metabólico relevantes na população adulta, de modo que a mensuração complementar da circunferência da cintura se faz relevante para uma classificação mais precisa e clinicamente correta. Na Subamostra Laboratorial da PNS 2013, a aferição da circunferência da cintura seguiu o protocolo padronizado pelo IBGE, com medição realizada no ponto médio entre a borda inferior da última costela e a borda superior da crista ilíaca, em conformidade com a recomendação metodológica da Organização Mundial da Saúde \cite{IBGE_PNS2013_Exames}.

O estudo longitudinal de \textcite{Son2016} evidencia que a razão Cintura-Estatura (RCEst), derivada da divisão entre a circunferência da cintura e a altura do indivíduo, deve ser incluída como variável auxiliar a fim de capturar a adiposidade visceral ajustada pela estatura individual, já que essa dimensão não é abrangida de forma integral pelo IMC ou pela circunferência da cintura em valor absoluto isoladamente. O estudo demonstra que a RCEst apresenta a vantagem de contextualizar o depósito de gordura central em relação ao porte corporal do indivíduo, de forma que dois indivíduos com a mesma circunferência absoluta possam apresentar perfis de risco visceral distintos em função da estatura. A partir dessa fundamentação, o bloco antropométrico candidato reúne três aspectos complementares do fenótipo corporal: a magnitude geral do acúmulo adiposo (IMC), a concentração de gordura na região abdominal (CC) e a adiposidade visceral ajustada pela estrutura corporal (RCEst) \cite{mlynarska2025type, Feller2010,Son2016}. A permanência de cada uma dessas variáveis no modelo final ficou condicionada ao diagnóstico de multicolinearidade descrito na Metodologia.

Segundo a \textcite{SBD2019}, a Hipertensão Arterial Sistêmica (HAS), em conjunto com a dislipidemia\footnote{Nome dado aos níveis alterados (geralmente altos) de gorduras, como colesterol e triglicerídeos, no sangue.} e a adiposidade central, forma o ``tripé'' clássico da Síndrome Metabólica, condição clínica de alta prevalência na população brasileira e fortemente associada à resistência à insulina e ao desenvolvimento do DM2. No presente trabalho, a pressão arterial foi considerada em duas dimensões complementares: a medida objetiva (pressões sistólica e diastólica aferidas durante a coleta da PNS 2013) e o autorrelato diagnóstico de HAS. Optou-se por considerar as duas fontes em razão da premissa central do trabalho, a qual prioriza dados objetivos sobre o autorrelato em razão do fenômeno documentado de subnotificação diagnóstica na população brasileira por \textcite{muzy2021prevalencia}, ao mesmo tempo em que preserva a dimensão de consciência diagnóstica do indivíduo como proxy\footnote{Variáveis utilizadas como indicadores indiretos de um conceito que não pode ser mensurado diretamente.} de acesso ao sistema de saúde.

A dislipidemia, operacionalizada por meio do diagnóstico autorreferido de colesterol alto (variável Q060 da PNS 2013), é incorporada ao modelo em razão do seu papel consolidado na fisiopatologia da Síndrome Metabólica \cite{mlynarska2025type, SBD2019}. Adicionalmente, a partir da combinação de HAS combinada, dislipidemia autorreferida e presença de obesidade (definida por IMC $\geq$ 30 kg/m² ou CC $\geq$ 80 cm em mulheres ou $\geq$ 94 cm em homens \cite{SBD2019}), foi desenvolvida a variável derivada \textit{Multimorbidade Metabólica}, operacionalizada como escore ordinal de zero a três, a fim de viabilizar a captura de interações sinérgicas entre condições clínicas que, isoladamente, poderiam não atingir o limiar de risco identificável pelo algoritmo, mas que em combinação elevam expressivamente a probabilidade de hiperglicemia crônica clinicamente significativa \cite{SBD2019, rzewuska2017epidemiology}.

Por fim, o estado de saúde autoavaliado (variável N001 da PNS 2013), internacionalmente referido como \textit{Self-Rated Health} (SRH), foi incorporado ao conjunto de classificadores em virtude da sua consolidação, em estudos epidemiológicos transversais e longitudinais, como preditor de desfechos crônicos e de mortalidade em populações adultas \cite{Rantanen2019}. No contexto específico da identificação de DM2 não diagnosticado, o estudo de \textcite{Rantanen2019} demonstra que a SRH apresenta uma alta relevância, já que, segundo os autores, indivíduos com hiperglicemia crônica não diagnosticada frequentemente manifestam sintomas genéricos em estágios pré-clínicos, a exemplo de fadiga, sede anormal leve e poliúria\footnote{Termo médico utilizado para descrever a produção e eliminação de quantidades excessivas de urina.}, capaz de impactar a percepção subjetiva de saúde antes mesmo da formalização diagnóstica.

\section{Mineração de Dados} \label{sec:mineracao}

De acordo com \textcite{mining2006introduction}, a mineração de dados (\textit{Data Mining}) configura-se como um conjunto de ferramentas e algoritmos que permitem a análise e solução de problemas que demandariam esforço manual massivo ou simplesmente ficariam sem solução. A mesma também pode ser definida como um conjunto de ferramentas ou técnicas que possibilitam identificar relações, padrões e regularidades que não seriam detectáveis por análise humana direta. No contexto do presente trabalho, a aplicação dessas técnicas sobre bases de dados populacionais, como os microdados da PNS 2013, viabiliza a transformação de informações clínicas e sociodemográficas coletadas em inteligência epidemiológica ativa, de maneira que possíveis padrões presentes na população possam ser identificados de forma sistemática \cite{islam2024comprehensive}.

A fim de orientar projetos de mineração de dados de forma estruturada e reprodutível, o modelo \textit{Cross Industry Standard Process for Data Mining} (CRISP-DM) foi desenvolvido em 1996 por uma parceria de empresas europeias liderada por Daimler-Benz, SPSS (\textit{Statistical Package for the Social Sciences}), NCR (\textit{National Cash Register}) e Ohra, sendo formalizado e publicado por Wirth e Hipp em 2000 \cite{wirth2000crisp}. O CRISP-DM organiza o ciclo de vida de um projeto de mineração de dados em seis fases iterativas: Compreensão do Negócio (\textit{Business Understanding}), Compreensão dos Dados (\textit{Data Understanding}), Preparação dos Dados (\textit{Data Preparation}), Modelagem (\textit{Modeling}), Avaliação (\textit{Evaluation}) e Implantação (\textit{Deployment}) \cite{wirth2000crisp}. A natureza iterativa do modelo é um dos seus diferenciais fundamentais, já que as fases não são executadas de forma linear e sequencial, e sim de maneira cíclica, de modo que ações e descobertas realizadas em fases avançadas do processo possam justificar retornos a etapas anteriores a fim de refinar os dados, as hipóteses ou as decisões de modelagem \cite{wirth2000crisp}. Diante disso, o CRISP-DM favorece a rastreabilidade metodológica desde a concepção inicial do problema até a validação da solução, característica relevante para a reprodutibilidade científica de projetos baseados em dados públicos.

A predominância do CRISP-DM como \textit{framework} para projetos de ciência de dados e mineração de dados é amplamente corroborada pela literatura especializada. Conforme evidenciado pela revisão sistemática conduzida por \textcite{schroer2021systematic}, o CRISP-DM consolida-se como o modelo mais adotado em projetos de \textit{Machine Learning} e mineração de dados ao redor do mundo, em virtude de sua flexibilidade para adaptação a diferentes domínios de aplicação e da clareza com que organiza as responsabilidades e entregas de cada fase do projeto. Embora a revisão aponte que o modelo original pode necessitar de adaptações de \textit{Design Thinking} para atender a requisitos modernos de interação com o usuário final, a robustez do CRISP-DM para o tratamento de dados estruturados consolida o mesmo como uma base ideal para o processamento de microdados populacionais como os que estão presentes na Subamostra Laboratorial da PNS 2013 \cite{schroer2021systematic}.

Conforme documentado por \textcite{shimaoka2024evolution}, o modelo original do CRISP-DM foi objeto de revisões e extensões metodológicas que buscaram alinhar suas fases ao ciclo de desenvolvimento de sistemas de inteligência artificial, incorporando práticas contemporâneas de engenharia de dados, gestão de experimentos e monitoramento de modelos em produção. Diante disso, essa metodologia permanece como referência central para projetos que, como o presente trabalho, combinam processamento de dados com desenvolvimento de modelos classificatórios e implantação de protótipos funcionais \cite{shimaoka2024evolution}.

\section{Pré-processamento e Tratamento de Dados} \label{sec:pre-processamento}

\textcite{hastie2009elements} abordam extensivamente como a qualidade dos dados submetidos ao treinamento de um modelo classificatório constitui um dos determinantes mais críticos de seu desempenho. Em bases de dados baseadas em inquéritos populacionais como a Subamostra Laboratorial da PNS 2013, é comum a presença de valores ausentes, decorrentes de não respostas, erros de coleta ou do desenho condicional do questionário, ponto abordado diretamente na nota técnica disponibilizada sobre a amostra \cite{IBGE}. Soma-se a isso o desbalanceamento entre as classes do desfecho, situação na qual a classe de interesse representa proporção significativamente menor que a classe majoritária \cite{pertiwi2020comparison}. Por esta razão, a etapa de pré-processamento dos dados se faz necessária para garantir a integridade estatística do conjunto de treino e a validade das associações aprendidas pelo algoritmo ao longo do processo de modelagem.

Além disso, a presença de valores ausentes é um dos desafios mais comuns no tratamento de dados \cite{hastie2009elements}. Dentre as abordagens disponíveis, a substituição por medidas de tendência central, como a substituição pela média ou pela mediana da variável, é uma alternativa de simples implementação, contudo apresenta uma limitação, já que ao desconsiderar as relações multivariadas entre os preditores, pode introduzir distorções na distribuição dos dados, produzindo registros implausíveis, onde a combinação de valores não corresponde a nenhum perfil real da população estudada \cite{hastie2009elements}. Já a imputação por meio do algoritmo \textit{K-Nearest Neighbors} (KNN) opera sobre a estrutura de similaridade multivariada dos dados, ou seja, cada valor ausente é preenchido com base na média ponderada dos $k$ vizinhos mais próximos no espaço multidimensional das demais variáveis disponíveis \cite{pedregosa2011scikit, pertiwi2020comparison}. Por meio desta abordagem, os valores substituídos preservam a estrutura de correlação entre si, de forma com que os registros resultantes sejam coerentes com os padrões observados na base de dados.

Em relação ao tratamento do desbalanceamento de classes, a literatura discute diversas técnicas. No contexto do presente trabalho, concentrou-se em dois tipos de métodos, a reamostragem no nível dos dados (\textit{data-level}), como o SMOTE (\textit{Synthetic Minority Over-sampling Technique}), e no nível algorítmico (\textit{algorithm-level}), como a aprendizagem sensível ao custo (\textit{Cost-Sensitive Learning}) \cite{araf2024costsensitive}. Conforme evidenciado pela revisão sistemática conduzida por \textcite{araf2024costsensitive}, as abordagens de reamostragem que geram instâncias sintéticas, como o SMOTE, apresentam limitações em contextos onde os dados originais são resultantes de desenhos amostrais complexos, visto que a criação de observações artificiais quebra a estrutura de ponderação amostral utilizada para representatividade populacional das estimativas \cite{10.1371/journal.pone.0280387}. 

Já a aprendizagem sensível ao custo opera diretamente sobre a função de perda do algoritmo, atribuindo penalizações aos erros de classificação de cada classe, fazendo com que o custo de classificar incorretamente um indivíduo da classe minoritária seja proporcionalmente superior ao custo de um erro na classe majoritária, sem a necessidade de alterar a distribuição original dos dados de treino \cite{araf2024costsensitive}. No contexto da biblioteca \textit{scikit-learn} \cite{pedregosa2011scikit}, essa estratégia é realizada por meio do hiperparâmetro \texttt{class\_weight='balanced'}, cuja função é atribuir pesos inversamente proporcionais à frequência de cada classe no conjunto de treino, ou seja, o algoritmo é induzido a tratar os erros de classificação da classe minoritária com maior peso durante a otimização \cite{pedregosa2011scikit}. Essa característica torna a aprendizagem sensível ao custo adequada a este trabalho, pois a preservação da estrutura amostral original viabiliza a manutenção dos pesos laboratoriais (\textit{Peso\_Lab}) recalculados pelo IBGE \cite{ibge2013pns}, de forma com que a representatividade populacional das estimativas produzidas na fase de Análise Exploratória de Dados não seja comprometida pelo tratamento do desbalanceamento na fase de modelagem.


\section{Algoritmos de \textit{Machine Learning}} \label{sec:algoritmos-ml}

A seleção dos algoritmos de classificação supervisionada constitui uma das decisões de maior impacto em projetos de \textit{Machine Learning}. No presente trabalho, foram selecionados dois algoritmos: a Regressão Logística, utilizada como modelo de linha de base (\textit{baseline}) interpretável, e o \textit{Random Forest}, escolhido como modelo de maior capacidade representacional para a captura de interações não lineares entre as variáveis classificadoras.

A Regressão Logística constitui um modelo estatístico clássico e linear cuja função é estimar a probabilidade de ocorrência de um desfecho binário a partir de um conjunto de variáveis explicativas \cite{hastie2009elements}. O modelo opera por meio da função logística, que transforma a combinação linear dos preditores em uma probabilidade restrita ao intervalo entre zero e um, viabilizando a classificação binária de indivíduos em relação ao desfecho de interesse \cite{hastie2009elements}. Conforme evidenciado por \textcite{nusinovici2020logistic}, a Regressão Logística apresenta desempenho comparável ao de algoritmos de \textit{Machine Learning} mais complexos em tarefas de classificação de doenças crônicas, de forma com que sua adoção como \textit{baseline}\footnote{Modelo inicial simples usado como ponto de partida para comparar e justificar o desempenho de algoritmos mais complexos.} se justifica tanto pela transparência paramétrica quanto pela consolidação histórica em estudos epidemiológicos. A interpretabilidade do modelo reside na possibilidade de extrair, a partir dos coeficientes estimados, a razão de chances (\textit{odds ratio}) associada a cada variável, métrica amplamente utilizada em epidemiologia clínica para quantificar a magnitude da associação entre um preditor e o desfecho \cite{nusinovici2020logistic}.

O \textit{Random Forest} constitui um algoritmo de aprendizagem por conjunto (\textit{ensemble}) que combina as classificações de múltiplas árvores de decisão independentes a fim de produzir uma classificação final mais robusta e generalizada do que qualquer árvore individual seria capaz de fornecer \cite{islam2024comprehensive}. Cada árvore é treinada sobre uma subamostra aleatória do conjunto de treino, e em cada nó de divisão apenas um subconjunto aleatório de variáveis é avaliado como candidato, de maneira em que as árvores individuais tornam-se descorrelacionadas entre si e a agregação das saídas das árvores produz classificações de menor variância \cite{islam2024comprehensive}. A principal vantagem do \textit{Random Forest} em relação à Regressão Logística, no contexto do presente trabalho, reside na sua capacidade de capturar interações não lineares entre os classificadores, a exemplo da interação entre idade avançada, circunferência da cintura elevada e baixa escolaridade em regiões de menor acesso à atenção primária, padrões que um modelo linear não conseguiria representar de forma adequada \cite{islam2024comprehensive}. Adicionalmente, a compatibilidade do \textit{Random Forest} com o método \textit{TreeExplainer} da biblioteca SHAP viabiliza a decomposição das contribuições individuais de cada variável para cada classificação realizada pelo modelo, requisito para a interpretabilidade clínica das saídas do protótipo \cite{salih2025perspective}.

\section{Avaliação de Modelos Classificatórios} \label{sec:avaliacao}

A avaliação do desempenho de modelos classificatórios constitui a etapa central no desenvolvimento de sistemas de \textit{Machine Learning} aplicados à saúde, visto que a escolha das métricas deve estar alinhada aos objetivos clínicos do instrumento e às consequências assimétricas dos diferentes tipos de erro classificatório \cite{Swaminathan2024}. Nesse sentido, a Matriz de Confusão é o instrumento central de avaliação em tarefas de classificação binária, organizando os quatro possíveis resultados de uma classificação: Verdadeiros Positivos (VP), Verdadeiros Negativos (VN), Falsos Positivos (FP) e Falsos Negativos (FN) \cite{Swaminathan2024}. A partir dessa estrutura, resultam as métricas utilizadas no presente trabalho.

A Sensibilidade (\textit{Recall}) expressa a capacidade do modelo de identificar corretamente os indivíduos portadores da condição, configurando-se como a métrica prioritária em instrumentos de triagem populacional:

\begin{equation}
\text{Sensibilidade} = \frac{VP}{VP + FN} \quad \text{ \cite{Swaminathan2024}}
\end{equation}

\noindent A Especificidade quantifica a capacidade do modelo de identificar corretamente os indivíduos saudáveis, de modo que valores altos implicam baixa taxa de encaminhamentos desnecessários:

\begin{equation}\text{Especificidade} = \frac{VN}{VN + FP} \text{ \cite{Swaminathan2024}}
\end{equation}

\noindent O \textit{F1-Score} constitui a média harmônica entre Precisão e Sensibilidade, cuja função é equilibrar ambas as métricas em cenários de desbalanceamento de classes \cite{Swaminathan2024}:

\begin{equation}F1 = 2 \times \frac{\text{Precisão} \times \text{Sensibilidade}}{\text{Precisão} + \text{Sensibilidade}} \text{ \cite{Swaminathan2024}}
\end{equation}

\noindent A Área Sob a Curva ROC (AUC-ROC) avalia a capacidade discriminativa global do modelo ao plotar a Sensibilidade em função da taxa de Falsos Positivos ao longo de todos os limiares de classificação possíveis. Valores próximos a 1 indicam elevada capacidade discriminativa, enquanto 0,5 equivale à classificação aleatória \cite{fawcett2006introduction, powers2020evaluation}.

A Validação Cruzada Estratificada com $k = 5$ \textit{folds} divide o conjunto de treino em $k$ partições, utilizando, em cada iteração, $k-1$ partições para treinamento e uma para validação \cite{pedregosa2011scikit}. A estratificação garante que a proporção da classe positiva seja preservada em cada partição, característica relevante em cenários de desbalanceamento \cite{pertiwi2020comparison}. A média e o desvio padrão das métricas ao longo das $k$ iterações descrevem a variação do desempenho entre partições.


\section{Inteligência Artificial Explicável (XAI)} \label{sec:xai}

A crescente adoção de algoritmos de aprendizagem de máquina em diversos contextos trouxe consigo um desafio: algoritmos de alta capacidade, como o \textit{Random Forest}, operam por meio de estruturas computacionais de alta complexidade, nas quais a classificação final resulta da agregação de centenas ou milhares de árvores de decisão individuais, de maneira em que a contribuição de cada variável para o resultado final torna-se vaga e praticamente impossível de interpretar por meio de análise direta \cite{salih2025perspective}. Esse fenômeno, amplamente referido na literatura como o dilema da ``caixa-preta'' (\textit{black box}), caracteriza-se como barreira à adoção de modelos de \textit{Machine Learning} em ambientes clínicos, visto que profissionais de saúde e gestores do sistema público necessitam compreender os fundamentos das classificações produzidas a fim de assumir a responsabilidade ética pelas condutas assistenciais delas derivadas \cite{salih2025perspective}.

No contexto da Atenção Primária à Saúde brasileira, essa barreira apresenta dimensões adicionais. A literatura evidencia que a resistência profissional à adoção de tecnologias digitais está entre os principais desafios à continuidade do modelo assistencial na APS \cite{geremia2020atenccao}, e a aceitabilidade de um sistema classificatório por parte das equipes multidisciplinares depende da capacidade do sistema de justificar suas saídas em linguagem compreensível ao usuário. Diante disso, a integração de técnicas de Inteligência Artificial Explicável (XAI) ao modelo proposto no presente trabalho não é apenas uma escolha técnica, mas um requisito de viabilidade do protótipo desenvolvido \cite{geremia2020atenccao, salih2025perspective}.

A Inteligência Artificial Explicável (XAI) é caracterizada como um conjunto de métodos, técnicas e ferramentas cuja função é transformar modelos de aprendizagem de máquina em sistemas onde as decisões sejam compreensíveis e auditáveis por usuários humanos, sejam eles especialistas técnicos ou profissionais de domínio \cite{salih2025perspective}. Dentre os métodos de XAI disponíveis na literatura, o SHAP (\textit{SHapley Additive exPlanations}) se sobressai como a abordagem de maior rigor matemático e mais ampla adoção em aplicações clínicas \cite{salih2025perspective}. O SHAP fundamenta-se na teoria dos jogos cooperativos, mais especificamente no conceito de valores de Shapley, desenvolvido por Lloyd Shapley em 1953, que define a contribuição de cada participante em um jogo cooperativo com base na sua contribuição marginal média em todas as coalizões possíveis de participantes \cite{salih2025perspective}. Trazendo este conceito para o contexto da explicabilidade de modelos de \textit{Machine Learning}, essa teoria permite que cada variável classificatória seja tratada como um ``jogador'' e que a classificação final do modelo seja decomposta na contribuição individual de cada variável.

A análise via SHAP abrange duas dimensões de interpretabilidade. A dimensão global, calculada sobre o conjunto de teste completo, produz uma visão agregada da importância relativa de cada variável nas classificações do modelo e permite confrontar as associações aprendidas pelo algoritmo com a evidência epidemiológica consolidada \cite{salih2025perspective}. A dimensão local decompõe a contribuição de cada variável no escore atribuído a um indivíduo, de modo que cada classificação venha acompanhada da justificativa do algoritmo \cite{salih2025perspective}. Os valores de Shapley expressam contribuições estatísticas observadas em dados transversais e não constituem relações causais nem predições de desenvolvimento futuro da condição clínica avaliada.

\section{Multicolinearidade e seus Efeitos sobre a Modelagem e a Interpretabilidade}\label{sec:multicolinearidade}

Segundo \textcite{Salih2024} a multicolinearidade é definida como um fenômeno no qual múltiplas características em modelos de aprendizado de máquina são altamente correlacionadas entre si. Conforme abordado por \textcite{Senaviratna2019}, no contexto da Regressão Logística, a multicolinearidade infla a variância dos coeficientes estimados, tornando as estimativas instáveis e comprometendo a confiabilidade dos testes de significância individual. No que se refere aos modelos baseados em árvores de decisão, como o \textit{Random Forest}, a multicolinearidade apresenta impacto reduzido sobre o desempenho classificatório, visto que esses algoritmos selecionam variáveis de forma iterativa com base em critério de ganho de informação, sendo naturalmente mais robustos à presença de correlações entre preditores \cite{Chan2022}.

Entretanto, \textcite{Salih2024} evidencia que a robustez classificatória dos modelos de árvore não se estende ao plano da interpretabilidade via SHAP, visto que, na versão original, o método assume independência estatística entre variáveis para a geração das explicações. Quando duas variáveis possuem forte correlação entre si e se correlacionam com o desfecho, apenas uma delas tende a receber um valor SHAP elevado, já que a inclusão da segunda variável no modelo não acrescenta ganho informacional significativo em relação à primeira. Assim, o valor SHAP resultante não captura o impacto real de cada variável sobre o desfecho quando as \textit{features} são colineares. Dessa forma, a lista de variáveis mais relevantes produzida pelo método pode se tornar incompleta ou enviesada em favor de uma das variáveis do par correlacionado \cite{Salih2024}. Esse fenômeno é particularmente relevante para o presente trabalho, visto que algumas das variáveis selecionadas, apresentadas na Metodologia, são clinicamente complementares, porém estatisticamente correlacionadas entre si. 

A revisão conduzida por \textcite{Salih2024} identificou sete abordagens propostas na literatura para mitigar o impacto da multicolinearidade sobre os métodos de XAI, dentre as quais se destacam o \textit{Modified Index Position} (MIP), o \textit{Extended Kernel SHAP} e o \textit{Multi-Collinearity Corrected} (MCC). Entretanto, o autor constata que nenhum método de XAI disponível mitiga, por natureza, o impacto da multicolinearidade sobre suas explicações, de forma com que o tratamento desse fenômeno demanda intervenções metodológicas complementares por parte do pesquisador \cite{Salih2024}.

